 \chapter{Actions, opérations sur le capital et indices}\label{ch:m1:shares-corporate-actions-indices}

Le vendredi 7 juin 2024, l'action Nvidia a clôturé à 1\,208.88~\$. Le lundi
matin, elle a ouvert autour de 121~\$. Personne n'a perdu un centime~: durant le week-end, chaque
détenteur avait reçu neuf actions supplémentaires pour chaque action détenue. Un graphique des
cours de clôture bruts fait apparaître un effondrement de 90~\%~; un backtest mené sur ce
graphique révèle une stratégie de vente à découvert remarquable~; un indice qui intégrait le
titre à son nouveau cours sans autre précaution aurait chuté avec lui.
Les actions ne sont pas des objets statiques. Elles distribuent des \omterm{def:m1:shares-corporate-actions-indices:dividend}{dividendes}, se divisent, se multiplient, fusionnent et
disparaissent, et tout système qui stocke un cours doit savoir ce qu'il est arrivé à
l'actif évalué. Ce chapitre traite de ces événements, de l'arithmétique
qui les annule, et des indices qui doivent être maintenus à travers
tous ces changements.

\section{Ce qu'est une action}

\begin{definition}[Action et capitalisation boursière]\label{def:m1:shares-corporate-actions-indices:share}
Une \emph{action}\index{action} est une part de propriété dans une société, conférant à son
détenteur un droit proportionnel sur ce que la société distribue et, en général, un
droit de vote. Avec $N$ actions en circulation au cours $P$, la
\emph{capitalisation boursière}\index{capitalisation boursière} est $NP$.
\end{definition}

\begin{definition}[Flottant]\label{def:m1:shares-corporate-actions-indices:float}
Le \emph{flottant}\index{flottant} d'une société est la fraction de ses
actions disponibles pour les investisseurs du public, à l'exclusion des participations
qui ne sont pas destinées à être cédées dans le cours normal des affaires~: fondateurs et autres
détenteurs stratégiques, États, participations croisées, actions bloquées après l'admission à la cote.
\end{definition}

Une société peut avoir plusieurs \emph{catégories} d'actions aux droits de vote ou
de \omterm{def:m1:shares-corporate-actions-indices:dividend}{dividendes} différents, chacune avec son propre cours et son propre symbole boursier~; la même action peut
aussi être cotée sur plusieurs bourses, dans plusieurs devises (voir le \cref{ch:m1:delta-one-instruments}).
(voir le \cref{ch:m1:delta-one-instruments}). «~Le cours de la société~» est
donc une convention~; les structures de données du \cref{ch:m1:reading-market-data}
existent pour la préciser.

\section{Dividendes}

\begin{definition}[Dividende, date d'enregistrement et date de détachement]\label{def:m1:shares-corporate-actions-indices:dividend}
Un \emph{dividende}\index{dividende} est une distribution de liquidités (ou d'actions) par une
société à ses actionnaires, d'un montant déclaré par action. Il est versé à
quiconque figure au registre à la clôture de la
\emph{date d'enregistrement}\index{date d'enregistrement}. La
\emph{date de détachement}\index{date de détachement} est le premier jour où un
acheteur de l'action n'a plus droit au dividende~: les transactions conclues le jour même ou
après celle-ci sont réglées trop tard pour figurer à temps sur le registre.
\end{definition}

\begin{dated}{2026-09}{Date de détachement et date d'enregistrement aux États-Unis}\label{dat:m1:shares-corporate-actions-indices:exdate}
Depuis que le règlement-livraison aux États-Unis est passé à un jour le 28 mai 2024, la \omterm{def:m1:shares-corporate-actions-indices:dividend}{date de détachement}
d'une distribution ordinaire est la \omterm{def:m1:shares-corporate-actions-indices:dividend}{date d'enregistrement} elle-même~: les règles des bourses
le précisent désormais. Avec un règlement-livraison à deux jours, la \omterm{def:m1:shares-corporate-actions-indices:dividend}{date de détachement} tombait un jour ouvré \emph{avant}
la \omterm{def:m1:shares-corporate-actions-indices:dividend}{date d'enregistrement}.
\end{dated}

\begin{omfigure}
\begin{tikzpicture}[font=\small]
\draw[-{Stealth[length=2.2mm]},thick,omExo] (-0.4,0) -- (15,0);
\foreach \x/\top/\bot in {0.8/Déclaration/{montant et dates annoncés},5.2/{Dernier jour\\cum-dividende}/{les acheteurs reçoivent encore le dividende},8.6/{Date de détachement}/{le cours ouvre plus bas d'environ $D$},12.6/{Paiement}/{les espèces parviennent aux détenteurs}}{
  \draw[omExo,thick] (\x,-0.12) -- (\x,0.12);
  \node[font=\small\bfseries,anchor=south,align=center] at (\x,0.2) {\top};
  \node[font=\footnotesize,omExo,anchor=north,align=center,text width=3.1cm] at (\x,-0.22) {\bot};
}
\draw[omThm,thick,decorate,decoration={brace,amplitude=4pt}] (8.6,1.0) -- (12.6,1.0)
  node[midway,above=5pt,font=\footnotesize,omThm] {le détenteur a droit à $D$~: une créance, et non un prix};
\end{tikzpicture}

\omcaption{Les dates d'un \omterm{def:m1:shares-corporate-actions-indices:dividend}{dividende}. Avec un règlement-livraison à un jour, la \omterm{def:m1:shares-corporate-actions-indices:dividend}{date de détachement}
coïncide avec la \omterm{def:m1:shares-corporate-actions-indices:dividend}{date d'enregistrement}. Entre la \omterm{def:m1:shares-corporate-actions-indices:dividend}{date de détachement} et le paiement, la
richesse du détenteur est l'action \emph{plus} le \omterm{def:m1:shares-corporate-actions-indices:dividend}{dividende} dû~: un système de P\&L
qui oublie la seconde partie affiche une perte qui ne s'est pas produite
(voir la \cref{met:m1:pnl-and-positions:explain})}
\end{omfigure}

\begin{proposition}[La baisse à la date de détachement]\label{prop:m1:shares-corporate-actions-indices:drop}
En négligeant les impôts et l'intérêt d'une nuit, et en l'absence de nouvelle entre les deux
enchères, une action qui clôture \omterm{def:m1:shares-corporate-actions-indices:dividend}{cum-dividende} au cours $P$ ouvre \omterm{def:m1:shares-corporate-actions-indices:dividend}{ex-dividende} à $P -
D$.
\end{proposition}
\begin{proof}
Une action achetée à la clôture \omterm{def:m1:shares-corporate-actions-indices:dividend}{cum-dividende} vaut, le lendemain matin, une action
\omterm{def:m1:shares-corporate-actions-indices:dividend}{ex-dividende} augmentée d'une créance de $D$ payable dans quelques jours. Si elle ouvrait au-dessus de $P - D$,
acheter à la clôture et vendre à l'ouverture rapporterait un profit sans risque~;
en dessous, l'opération inverse en produirait un.
\end{proof}

En pratique, la baisse est un peu inférieure à $D$ lorsque les \omterm{def:m1:shares-corporate-actions-indices:dividend}{dividendes} sont imposés plus
lourdement que les plus-values, car le détenteur marginal valorise un dollar de
\omterm{def:m1:shares-corporate-actions-indices:dividend}{dividende} à moins d'un dollar~; l'écart, ainsi que le prêt d'actions
à cheval sur la \omterm{def:m1:shares-corporate-actions-indices:dividend}{date de détachement} qui l'exploite, sont traités au
\cref{ch:m1:stock-loan-and-short-selling}.

\section{Divisions d'actions, augmentations de capital avec droits et scissions}

\begin{definition}[Division d'actions]\label{def:m1:shares-corporate-actions-indices:split}
Dans une \emph{division d'actions}\index{division d'actions} de $k$ pour 1, chaque action devient $k$
actions. Le nombre d'actions est multiplié par $k$, le cours est divisé par
$k$, et rien d'autre ne change. Un \emph{regroupement d'actions} a $k<1$.
\end{definition}

\begin{definition}[Augmentation de capital avec droits]\label{def:m1:shares-corporate-actions-indices:rights}
Dans une \emph{augmentation de capital avec droits}\index{augmentation de capital avec droits}, une société lève des fonds en
donnant à ses actionnaires le droit de souscrire $n$ actions nouvelles pour chaque lot de $N$ actions détenues,
à un prix de souscription $S$ inférieur au cours du marché. Les droits sont en général
négociables pendant quelques semaines.
\end{definition}

\begin{proposition}[Cours théorique ex-droits]\label{prop:m1:shares-corporate-actions-indices:terp}
Si l'action clôture à $P$ avant une \omterm{def:m1:shares-corporate-actions-indices:rights}{augmentation de capital avec droits} $n$ pour $N$ au prix $S$, son
cours théorique une fois les droits détachés est
\[
\mathrm{TERP} = \frac{NP + nS}{N+n},
\],
et un droit de souscription à une action nouvelle vaut $\mathrm{TERP} - S$.
\end{proposition}
\begin{proof}
Après l'augmentation, $N+n$ actions détiennent l'ancienne société, valant $NP$, plus $nS$ de
liquidités nouvelles. Un droit permet à son détenteur de payer $S$ pour un actif valant
$\mathrm{TERP}$.
\end{proof}

\begin{example}[Un pour quatre à six]\label{ex:m1:shares-corporate-actions-indices:terp}
Une action se négocie à 10.00~; la société émet 1 action nouvelle pour 4 détenues, à 6.00.
$\mathrm{TERP} = (40 + 6)/5 = 9.20$. Un droit donnant accès à une action nouvelle vaut
$3.20$~; le détenteur de 4 actions reçoit un droit et possède $4 \times 9.20 +
3.20 = 40.00$, comme auparavant. Le cours de l'action baisse de 8~\% à la \omterm{def:m1:shares-corporate-actions-indices:dividend}{date de détachement},
et personne n'est plus pauvre~; le détenteur qui ignore les droits, en revanche, abandonne 3.20.
\end{example}

\begin{definition}[Scission]\label{def:m1:shares-corporate-actions-indices:spinoff}
Dans une \emph{scission}\index{scission} une société distribue à ses
actionnaires les actions d'une filiale, qui devient une société
cotée séparément. Le cours de la société mère baisse à la \omterm{def:m1:shares-corporate-actions-indices:dividend}{date de détachement}
d'environ la valeur distribuée par action.
\end{definition}

Les fusions complètent la liste~: dans une fusion en numéraire, l'action disparaît contre un
paiement~; dans une fusion en actions, elle devient un nombre fixe d'actions de l'acquéreur.
Pour les données, les deux mettent fin à une série de cours, et la règle est la même
que partout dans ce chapitre~: la \emph{richesse du détenteur} est continue~; le
\emph{cours} ne l'est pas.

\section{Facteurs d'ajustement}

\begin{definition}[Facteur d'ajustement]\label{def:m1:shares-corporate-actions-indices:factor}
Le \emph{facteur d'ajustement}\index{facteur d'ajustement} d'une opération sur titres
est le nombre $f$ par lequel on multiplie tout cours antérieur à sa \omterm{def:m1:shares-corporate-actions-indices:dividend}{date de détachement}
pour que la série ajustée ne présente aucun saut dû à l'opération. Un
\emph{cours rétro-ajusté} est le cours brut multiplié par le produit des facteurs
de toutes les opérations postérieures.
\end{definition}

\begin{method}[Le facteur de chaque opération]\label{met:m1:shares-corporate-actions-indices:factors}
Avec $P$ le dernier cours cum-droit~:
\begin{enumerate}
\item Division de $k$ pour 1~: $f = 1/k$ (les volumes sont multipliés par $k$).
\item \omterm{def:m1:shares-corporate-actions-indices:dividend}{Dividende} en numéraire $D$~: $f = (P-D)/P$.
\item \omterm{def:m1:shares-corporate-actions-indices:rights}{Augmentation de capital avec droits}~: $f = \mathrm{TERP}/P$.
\item \omterm{def:m1:shares-corporate-actions-indices:spinoff}{Scission} valant $V$ par action de la société mère~: $f = (P-V)/P$.
\end{enumerate}
Conserver les cours bruts et les opérations~; calculer les séries ajustées à la demande.
Les cours ajustés changent à chaque nouvelle opération, si bien qu'un cours ajusté
enregistré est un nombre dont la date n'est pas indiquée.
\end{method}

\begin{proposition}[Ce que fait exactement le facteur de dividende]\label{prop:m1:shares-corporate-actions-indices:exact}
Soit une action qui clôture cum-droit à $P$ et ex-droit à $P_{\mathrm{ex}}$. Le rendement de la
série ajustée par $f = (P-D)/P$ vaut $P_{\mathrm{ex}}/(P-D) - 1$, tandis que le
rendement total du détenteur vaut $(P_{\mathrm{ex}}+D)/P - 1$. Ils sont égaux lorsque
$P_{\mathrm{ex}} = P - D$ et diffèrent au second ordre sinon. Le facteur
$f^\star = P_{\mathrm{ex}}/(P_{\mathrm{ex}}+D)$, qui suppose que le \omterm{def:m1:shares-corporate-actions-indices:dividend}{dividende} est
réinvesti au cours ex, reproduit le rendement total dans tous les cas.
\end{proposition}
\begin{proof}
Le cours cum-droit ajusté vaut $fP = P - D$, ce qui donne le premier rendement. Avec
$f^\star$, le cours cum-droit ajusté vaut $P P_{\mathrm{ex}}/(P_{\mathrm{ex}}+D)$ et
le rendement vaut $(P_{\mathrm{ex}}+D)/P - 1$. Pour la différence, écrivons
$P_{\mathrm{ex}} = (P-D)(1+\varepsilon)$~: les deux rendements valent $\varepsilon$
et $\varepsilon\,(P-D)/P$, qui diffèrent de $\varepsilon D/P$.
\end{proof}

\begin{example}[Deux points de base]\label{ex:m1:shares-corporate-actions-indices:approx}
$P = 102$, $D = 2$, et le titre clôture ex-droit à 99, soit une baisse de 1~\% nette du
\omterm{def:m1:shares-corporate-actions-indices:dividend}{dividende}. Le facteur standard donne $-1.000\%$~; le détenteur a gagné $101/102 -
1 = -0.980\%$. L'écart, deux points de base, vaut $\varepsilon D/P$. Les fournisseurs
diffèrent par le facteur qu'ils publient~: deux séries «~clôture ajustée~» pour le
même titre ne sont pas forcément identiques, et une base de données de recherche doit
préciser sa convention.
\end{example}

\begin{omfigure}
\begin{tikzpicture}
\begin{axis}[width=11cm,height=5.2cm,ymode=log,
  xlabel={jour de négociation},ylabel={cours (échelle logarithmique)},
  tick label style={font=\small},label style={font=\small},
  xmin=0,xmax=250,ymin=25,ymax=900,grid=major,grid style={gray!25},
  ytick={30,50,100,200,400,800},yticklabels={30,50,100,200,400,800},
  legend columns=2,legend style={font=\footnotesize,at={(0.5,-0.3)},anchor=north,draw=none,fill=none,/tikz/every even column/.append style={column sep=8pt}},
  table/col sep=comma]
\addplot[omThm,thick] table[x=day,y=raw]{figdata/markets-1/08-shares-corporate-actions-indices/split.csv};
\addplot[omDef,very thick] table[x=day,y=adjusted]{figdata/markets-1/08-shares-corporate-actions-indices/split.csv};
\legend{cours brut de clôture,cours de clôture rétro-ajusté}
\end{axis}
\end{tikzpicture}

\omcaption{Une action simulée avec une division de dix pour un le jour 150 et des
\omterm{def:m1:shares-corporate-actions-indices:dividend}{dividendes} les jours 60 et 200. La série brute montre une chute de 90~\% que nul
détenteur n'a subie~; la série rétro-ajustée est continue, et \emph{toutes}
ses valeurs antérieures au jour 200 diffèrent des cours réellement négociés. Données~:
le script du chapitre.}\label{fig:m1:shares-corporate-actions-indices:split}
\end{omfigure}

\begin{remark}[Trois usages, trois séries]\label{rem:m1:shares-corporate-actions-indices:uses}
Utiliser les cours \emph{bruts} pour tout ce qui touche un ordre~: prix limites, pas de cotation,
exécutions, frais, simulateur. Utiliser les cours \emph{ajustés} pour les rendements, les
volatilités et les signaux calculés à cheval sur une \omterm{def:m1:shares-corporate-actions-indices:dividend}{date de détachement}. Et dans un
backtest, utiliser la série ajustée \emph{telle qu'elle était connue à chaque date}~: un signal
calculé aujourd'hui à partir des cours rétro-ajustés de 2019 a vu la division de 2024. Cette
discipline, appelée données point-in-time (telles que connues à chaque date), fait l'objet du livre~7.
\end{remark}

\section{Indices}

Un indice est un portefeuille défini par une règle, dont la valeur est publiée sous la forme
d'un seul nombre. Deux choix de conception le définissent~: la manière dont les composantes
sont pondérées, et la manière dont le nombre reste continu lorsque le portefeuille change.

\begin{definition}[Diviseur d'indice]\label{def:m1:shares-corporate-actions-indices:divisor}
Un indice pondéré par la capitalisation, dont les composantes $i$ ont pour cours $p_i$, pour
nombre d'actions en circulation $s_i$ et pour facteur de \omterm{def:m1:shares-corporate-actions-indices:float}{flottant} $\phi_i \in [0,1]$, a pour niveau
\[
I = \frac{1}{\Delta}\sum_i p_i\,s_i\,\phi_i ,
\]
où le \emph{diviseur d'indice}\index{diviseur d'indice} $\Delta$ est choisi à la création pour
donner un niveau de départ rond, et n'est ensuite modifié que pour compenser les variations
du numérateur qui ne sont pas des mouvements de cours.
\end{definition}

\begin{proposition}[Ajustement par le diviseur]\label{prop:m1:shares-corporate-actions-indices:divisoradj}
Lorsque, à des cours donnés, un changement de composantes, d'actions ou de \omterm{def:m1:shares-corporate-actions-indices:float}{flottant} fait passer la
valeur de marché de l'indice de $M$ à $M'$, le niveau est inchangé si et seulement si
le diviseur devient $\Delta' = \Delta\,M'/M$.
\end{proposition}
\begin{proof}
$I = M/\Delta = M'/\Delta'$.
\end{proof}

Un fractionnement modifie $p_i$ et $s_i$ dans des proportions opposées et laisse $M$
inchangé~: un indice pondéré par la capitalisation n'a pas besoin d'ajustement pour lui. Un
\emph{indice pondéré par les cours}, $I = \sum_i p_i/\Delta$, en a besoin~:
le poids de la composante fractionnée est divisé par dix, et le diviseur doit diminuer pour
compenser. Le plus connu des indices pondérés par les cours, le Dow Jones Industrial
Average, a été ajusté si souvent que son diviseur est nettement inférieur
à un~: une variation d'un dollar sur l'une quelconque de ses trente actions déplace
l'indice de plusieurs points.

\begin{definition}[Indice de rendement total]\label{def:m1:shares-corporate-actions-indices:tri}
Un \emph{indice de cours} ne reflète que les variations de cours~: à la \omterm{def:m1:shares-corporate-actions-indices:dividend}{date de détachement}, il recule
avec sa composante. Un \emph{indice de rendement total}\index{indice de rendement total}
réinvestit chaque \omterm{def:m1:shares-corporate-actions-indices:dividend}{dividende} dans l'indice à sa \omterm{def:m1:shares-corporate-actions-indices:dividend}{date de détachement} (un indice de rendement total \emph{net}, après déduction d'une
retenue à la source), et constitue le bon indice de référence pour
un fonds qui perçoit les \omterm{def:m1:shares-corporate-actions-indices:dividend}{dividendes}.
\end{definition}

\begin{omfigure}
\begin{tikzpicture}
\begin{axis}[width=11cm,height=4.8cm,ybar,bar width=7pt,
  ylabel={poids (\%)},symbolic x coords={A,B,C,D,E},xtick=data,
  tick label style={font=\small},label style={font=\small},
  ymin=0,ymax=70,ymajorgrids,grid style={gray!25},enlarge x limits=0.15,area legend,
  legend columns=3,legend style={font=\footnotesize,at={(0.5,-0.2)},anchor=north,draw=none,fill=none,/tikz/every even column/.append style={column sep=8pt}},
  table/col sep=comma]
\addplot[fill=omDef!70,draw=omDef] table[x=stock,y=cap]{figdata/markets-1/08-shares-corporate-actions-indices/weights.csv};
\addplot[fill=omThm!55,draw=omThm] table[x=stock,y=price]{figdata/markets-1/08-shares-corporate-actions-indices/weights.csv};
\addplot[fill=omMeth!55,draw=omMeth] table[x=stock,y=equal]{figdata/markets-1/08-shares-corporate-actions-indices/weights.csv};
\legend{capitalisation flottante ajustée,cours,égalitaire}
\end{axis}
\end{tikzpicture}

\omcaption{Cinq actions illustratives selon trois règles de pondération. L'action D, la plus petite
entreprise, domine l'indice pondéré par les cours parce que son cours est
de 900~; l'action B, l'une des plus grandes entreprises, y disparaît presque. Données~: le script du chapitre.}
\end{omfigure}

\begin{omfigure}
\begin{tikzpicture}
\begin{axis}[width=11cm,height=4.6cm,
  xlabel={jour de séance},ylabel={niveau de l'indice},
  tick label style={font=\small},label style={font=\small},
  xmin=0,xmax=59,ymin=900,ymax=1400,grid=major,grid style={gray!25},
  yticklabel style={/pgf/number format/1000 sep={\,}},
  legend columns=2,legend style={font=\footnotesize,at={(0.5,-0.34)},anchor=north,draw=none,fill=none,/tikz/every even column/.append style={column sep=8pt}},
  table/col sep=comma]
\addplot[omDef,very thick] table[x=day,y=maintained]{figdata/markets-1/08-shares-corporate-actions-indices/rebase.csv};
\addplot[omThm,thick,dashed] table[x=day,y=naive]{figdata/markets-1/08-shares-corporate-actions-indices/rebase.csv};
\legend{diviseur ajusté le jour 30,diviseur laissé inchangé}
\end{axis}
\end{tikzpicture}

\omcaption{Le jour 30, l'action B quitte un indice de cinq titres, et une entreprise plus grande y entre.
Avec la \cref{prop:m1:shares-corporate-actions-indices:divisoradj}, le niveau est continu~; sans elle, l'indice
«~gagne~» un tiers en une journée où
aucun cours n'a bougé. Données~: le script du chapitre.}
\end{omfigure}

\section{Tutoriel~: ajuster une série et maintenir un indice}\label{tut:m1:shares-corporate-actions-indices:adjust}

\begin{tutorial}
\textbf{Objectif.} Ajuster rétroactivement une série de cours à travers un fractionnement et deux \omterm{def:m1:shares-corporate-actions-indices:dividend}{dividendes},
puis faire évoluer un indice lors d'un changement de composante. \textbf{État final~:}
la \cref{fig:m1:shares-corporate-actions-indices:split} et les deux trajectoires d'indice
de la dernière figure.
\begin{tutsteps}
\item \textbf{Facteurs et ajustement rétroactif.} Chaque opération multiplie chaque
cours \emph{antérieur} par son facteur.
\omcode{code/markets-1/08-shares-corporate-actions-indices/python/corpact.py}{7}{14}{Facteurs de fractionnement et de dividende.}\label{lst:m1:shares-corporate-actions-indices:factors}
\omcode{code/markets-1/08-shares-corporate-actions-indices/python/corpact.py}{31}{37}{Ajustement rétroactif~: une ligne par opération.}\label{lst:m1:shares-corporate-actions-indices:backadj}
\item \textbf{Vérification.} Le script de la figure construit la série brute à partir d'une
trajectoire de rendement total connue et vérifie que la série ajustée a exactement
les rendements de cette trajectoire.
\item \textbf{L'indice.} Le diviseur ne change qu'à l'intérieur de \texttt{rebase}, et
seulement du rapport des valeurs de marché à cours inchangés.
\omcode{code/markets-1/08-shares-corporate-actions-indices/python/corpact.py}{41}{58}{Indice pondéré par la capitalisation flottante, avec maintien du diviseur.}\label{lst:m1:shares-corporate-actions-indices:index}
\item \textbf{Le casser.} Remplacer la composante sans appeler
\texttt{rebase}~: le niveau passe de 981 à 1\,320.
\end{tutsteps}
\textbf{Pour aller plus loin.} Ajouter une version en rendement total de l'indice~: à chaque
\omterm{def:m1:shares-corporate-actions-indices:dividend}{date de détachement}, augmenter le niveau des points d'indice du \omterm{def:m1:shares-corporate-actions-indices:dividend}{dividende}. Appliquer ensuite le
fractionnement de dix pour un sur une version pondérée par les cours, et calculer le nouveau
diviseur à la main.
\end{tutorial}

\section{Construction~: l'ajusteur d'opérations sur titres}\label{bld:m1:shares-corporate-actions-indices:adjuster}

\begin{build}
\bfield{Rôle} Un service de la petite entreprise répond à la question «~quel est le cours ajusté de ce
symbole à cette date, tel qu'il est connu à cette date~?~».
La recherche, le risque et le tenancier des P\&L l'appellent~; aucun ne stocke de cours
ajustés. \bfield{Interface} \texttt{Action(symbol, ex\_date, kind, params,
announced)} avec \texttt{kind} valant \texttt{split}, \texttt{cash\_dividend},
\texttt{rights}, \texttt{spinoff}; \texttt{Adjuster.add(action)};
\texttt{factor(symbol, date, as\_of)} renvoyant le produit des facteurs des
opérations avec \texttt{date < ex\_date <= as\_of}~; \texttt{adjust\_position(
position, action)} pour le tenancier du \cref{ch:m1:pnl-and-positions}.
\bfield{Règles} Les facteurs suivent
la \cref{met:m1:shares-corporate-actions-indices:factors}~; la convention de \omterm{def:m1:shares-corporate-actions-indices:dividend}{dividende}
est un argument du constructeur et figure dans chaque sortie. Une
opération est invisible pour les requêtes dont \texttt{as\_of} précède son annonce.
Un fractionnement multiplie la quantité d'une position par $k$ et divise son
\omterm{def:m1:pnl-and-positions:realised}{coût moyen} par $k$, sans toucher à la trésorerie~; un \omterm{def:m1:shares-corporate-actions-indices:dividend}{dividende} en numéraire
inscrit $qD$ au grand livre à la \omterm{def:m1:shares-corporate-actions-indices:dividend}{date de détachement}, en créance.
\bfield{Tests d'acceptation} \texttt{code/firm/corpactions/tests/}~: continuité
lors d'un fractionnement~; la règle de l'as-of~; l'invariance de la position (quantité
multipliée par le \omterm{def:m1:pnl-and-positions:realised}{coût moyen} inchangée par un fractionnement)~; l'exemple de l'\omterm{def:m1:shares-corporate-actions-indices:rights}{augmentation de capital avec droits}, une pour quatre.
\bfield{Pour aller plus loin} Fusions~: transformer une position dans la cible en trésorerie ou en
actions de l'acquéreur, puis clore la série de cours.
\end{build}

\begin{omsources}
\item CNBC, «~Nvidia announces 10-for-1 stock split~», 22 mai 2024; Yahoo
Finance, «~Nvidia stock rises after 10-for-1 stock split~», 10 juin 2024.
\item US Securities and Exchange Commission, Release 34-99881 (NYSE Arca
modification de règle sur les dates de détachement dans le cycle de règlement-livraison à un jour), 2024.
\item S\&P Dow Jones Indices, \emph{Index Mathematics Methodology}, avril
2026; \emph{Dow Jones Averages Methodology}, mai 2026; \emph{How the Dow
Works} (document pédagogique).
\item E.~Elton and M.~Gruber, «~Marginal stockholder tax rates and the
clientele effect~», \emph{Review of Economics and Statistics} 52 (1970) ---
la référence classique sur la baisse à la date de détachement.
\end{omsources}

\section{Exercices}

\begin{exercise}[$\star$]\label{exo:m1:shares-corporate-actions-indices:1}
Une société compte 800 millions d'actions à 65~\$ et un \omterm{def:m1:shares-corporate-actions-indices:float}{flottant} de 70~\%. Calculer sa
\omterm{def:m1:shares-corporate-actions-indices:share}{capitalisation boursière} et sa capitalisation ajustée du \omterm{def:m1:shares-corporate-actions-indices:float}{flottant}.
\end{exercise}

\begin{exercise}[$\star$]\label{exo:m1:shares-corporate-actions-indices:2}
Une action clôture à 240~\$ avant une \omterm{def:m1:shares-corporate-actions-indices:split}{division d'actions} 3-pour-2. Donner le cours
théorique d'ouverture, le \omterm{def:m1:shares-corporate-actions-indices:factor}{facteur d'ajustement} et la nouvelle position d'un détenteur de 500
actions.
\end{exercise}

\begin{exercise}[$\star$]\label{exo:m1:shares-corporate-actions-indices:3}
Une action clôture \omterm{def:m1:shares-corporate-actions-indices:dividend}{cum-dividende} à 80.00~\$ avec un \omterm{def:m1:shares-corporate-actions-indices:dividend}{dividende} de 1.20~\$. Donner le
cours théorique ex, le \omterm{def:m1:shares-corporate-actions-indices:factor}{facteur d'ajustement} et la valeur ajustée rétroactivement d'un cours de
76.00~\$ observé un mois plus tôt.
\end{exercise}

\begin{exercise}[$\star\star$]\label{exo:m1:shares-corporate-actions-indices:4}
Une société cotée à 25~\euro annonce une \omterm{def:m1:shares-corporate-actions-indices:rights}{augmentation de capital avec droits} 2-pour-7 à 16~\euro. Calculer
le TERP, la valeur d'un droit, le \omterm{def:m1:shares-corporate-actions-indices:factor}{facteur d'ajustement}, et vérifier qu'un détenteur de
700 actions qui vend tous ses droits reste aussi riche qu'auparavant.
\end{exercise}

\begin{exercise}[$\star\star$]\label{exo:m1:shares-corporate-actions-indices:5}
Un indice pondéré par les prix de trois actions à 50~\$, 120~\$ et 330~\$ a un diviseur
de 0.5. Donner son niveau. La troisième action fait l'objet d'une \omterm{def:m1:shares-corporate-actions-indices:split}{division d'actions} 3-pour-1~: donner le nouveau diviseur
et la nouvelle pondération de chaque action.
\end{exercise}

\begin{exercise}[$\star\star$]\label{exo:m1:shares-corporate-actions-indices:6}
Une action clôture \omterm{def:m1:shares-corporate-actions-indices:dividend}{cum-dividende} à 50.00~\$ avec un \omterm{def:m1:shares-corporate-actions-indices:dividend}{dividende} de 2.50~\$ et clôture à la
\omterm{def:m1:shares-corporate-actions-indices:dividend}{date de détachement} à 48.45~\$. Calculer le rendement à la \omterm{def:m1:shares-corporate-actions-indices:dividend}{date de détachement} avec le facteur standard, le
rendement total du détenteur, et l'écart. Vérifier cela avec la formule
de la \cref{prop:m1:shares-corporate-actions-indices:exact}.
\end{exercise}

\begin{exercise}[$\star\star\star$]\label{exo:m1:shares-corporate-actions-indices:7}
\emph{Programmation.} Un indice pondéré par la capitalisation compte deux actions~: X, avec 100
millions d'actions à 20~\$ et Y, avec 50 millions à 60~\$, des \omterm{def:m1:shares-corporate-actions-indices:float}{flottants} de 1, et un
niveau initial de 1\,000. En utilisant \texttt{CapIndex}~: (a) donner le diviseur~; (b) Y
rachète 10 millions d'actions à cours inchangés~: donner le nouveau diviseur~; (c)
les cours passent ensuite à 22~\$ et 57~\$~: donner le niveau.
\end{exercise}

\begin{exercise}[$\star\star\star$]\label{exo:m1:shares-corporate-actions-indices:8}
\emph{Trouver la faille.} Un chercheur télécharge aujourd'hui les cours de clôture quotidiens ajustés rétroactivement
pour 3\,000 actions sur vingt ans et teste la règle «~acheter les actions
dont le cours est inférieur à 5~\$~». Le backtest est excellent. Identifier deux manières distinctes
par lesquelles les cours ajustés ont contaminé le test.
\end{exercise}

\section{Problème~: Maintenir un indice pendant une mauvaise semaine}

\begin{problem}[{Problème du week-end --- cinq jours chez un fournisseur d'indices}]\label{pb:m1:shares-corporate-actions-indices:1}
On maintient un indice pondéré par la capitalisation ajustée du \omterm{def:m1:shares-corporate-actions-indices:float}{flottant} de quatre actions.
À la clôture du vendredi~:
\begin{center}\small
\begin{tabular}{@{}lrrr@{}}
\hline
Action & Cours & Actions (millions) & \omterm{def:m1:shares-corporate-actions-indices:float}{Flottant} \\
\hline
P & 40.00 & 500 & 1.00 \\
Q & 150.00 & 200 & 0.80 \\
R & 12.00 & 1\,000 & 0.50 \\
S & 75.00 & 400 & 0.90 \\
\hline
\end{tabular}
\end{center}
L'indice a clôturé à 2\,500.00.

\medskip\noindent\textbf{Partie I --- Vendredi.}
\begin{enumerate}
\item Calculer la capitalisation ajustée du \omterm{def:m1:shares-corporate-actions-indices:float}{flottant} de chaque action et la
valeur de marché de l'indice.
\item Donner le diviseur.
\item Donner la pondération de chaque action.
\item De combien de points d'indice l'indice varie-t-il si Q monte de 1~\%~?
\end{enumerate}

\medskip\noindent\textbf{Partie II --- Lundi~: une \omterm{def:m1:shares-corporate-actions-indices:split}{division d'actions} et un \omterm{def:m1:shares-corporate-actions-indices:dividend}{dividende}.}
\begin{enumerate}[resume]
\item Q fait l'objet d'une \omterm{def:m1:shares-corporate-actions-indices:split}{division d'actions} 3-pour-1 avant l'ouverture. Que deviennent son cours,
ses actions et le diviseur~?
\item P passe \omterm{def:m1:shares-corporate-actions-indices:dividend}{ex-dividende} pour un \omterm{def:m1:shares-corporate-actions-indices:dividend}{dividende} de 1.00 et ouvre à son cours théorique ex~;
rien d'autre ne bouge. Donner l'indice de cours. Le diviseur est-il ajusté~?
\item Donner le niveau de la version \emph{rendement total} de l'indice, qui
a aussi clôturé le vendredi à 2\,500.00.
\item Un fonds répliquant l'indice de cours a reçu le \omterm{def:m1:shares-corporate-actions-indices:dividend}{dividende}. Avec quel indice
doit-il être comparé~?
\end{enumerate}

\medskip\noindent\textbf{Partie III --- Mercredi~: une \omterm{def:m1:shares-corporate-actions-indices:rights}{augmentation de capital avec droits}.}
Les cours sont ceux de l'ouverture du lundi.
\begin{enumerate}[resume]
\item R annonce une \omterm{def:m1:shares-corporate-actions-indices:rights}{augmentation de capital avec droits} 1-pour-4 à 8.00, avec effet le mercredi. Donner
le TERP et le nouveau nombre d'actions de R.
\item En supposant l'augmentation entièrement souscrite, donner la nouvelle
capitalisation ajustée du \omterm{def:m1:shares-corporate-actions-indices:float}{flottant} de R au TERP et la nouvelle valeur de marché de l'indice.
\item Donner le nouveau diviseur.
\item Un fonds de réplication détient R à sa pondération dans l'indice. Que doit-il faire, en mots,
pour rester aligné sur l'indice~?
\end{enumerate}

\medskip\noindent\textbf{Partie IV --- Vendredi~: un remplacement.}
Les cours sont inchangés depuis le mercredi.
\begin{enumerate}[resume]
\item S est rachetée contre espèces et sort de l'indice à 75.00. Elle est remplacée
par T~: cours 30.00, 600~millions d'actions, \omterm{def:m1:shares-corporate-actions-indices:float}{flottant} 0.75. Donner la nouvelle
valeur de marché et le nouveau diviseur.
\item Donner les nouvelles pondérations.
\item Un fonds de 2~milliards de dollars qui réplique l'indice doit vendre ses actions S
et acheter des actions T. Combien de dollars représente chacune de ces opérations~?
\item T s'échange pour 150~millions de dollars par jour. Exprimer l'achat en jours
de volume et estimer son coût avec la \cref{met:m1:the-sell-side:sqrt} pour une
volatilité quotidienne de 2~\%.
\item Qui se trouve de l'autre côté de cet achat, et quand ces investisseurs ont-ils acheté~?
\item Le fournisseur d'indice a annoncé le changement cinq jours plus tôt. Pourquoi ne pas
annoncer et appliquer la modification le même jour~?
\item Énoncer le \emph{résultat nommé}~: le diviseur de l'indice à la fin de la
semaine.
\item En une phrase~: quelle est la seule chose qu'un ajustement du diviseur ne doit
jamais faire~?
\end{enumerate}
\end{problem}

\section{Questions d'entretien}

\begin{interviewq}[$\star$ \iqroles{chercheur, développeur, trader}]\label{iq:m1:shares-corporate-actions-indices:1}
Le cours d'une action chute de 50~\% du jour au lendemain, alors que la valeur de la société reste inchangée.
Que s'est-il passé, et que doit faire votre chaîne de traitement des données~?
\end{interviewq}

\begin{interviewq}[$\star$ \iqroles{trader, chercheur}]\label{iq:m1:shares-corporate-actions-indices:2}
De combien une action doit-elle baisser à sa \omterm{def:m1:shares-corporate-actions-indices:dividend}{date de détachement}~? Pourquoi pourrait-elle
chuter de moins~?
\end{interviewq}

\begin{interviewq}[$\star\star$ \iqroles{chercheur, ingénieur ML}]\label{iq:m1:shares-corporate-actions-indices:3}
Quand faut-il utiliser les cours bruts et quand les cours ajustés~? Donnez un bug causé
par chacun des mauvais choix.
\end{interviewq}

\begin{interviewq}[$\star\star$ \iqroles{trader, chercheur}]\label{iq:m1:shares-corporate-actions-indices:4}
Pourquoi une \omterm{def:m1:shares-corporate-actions-indices:split}{division d'actions} modifie-t-elle les pondérations d'un indice pondéré par les cours et non
celles d'un indice pondéré par la capitalisation~?
\end{interviewq}

\begin{interviewq}[$\star\star$ \iqroles{développeur}]\label{iq:m1:shares-corporate-actions-indices:5}
Concevez le stockage des cours et des opérations sur titres pour une base de données de recherche.
Que stocke-t-on, que calcule-t-on à la lecture, et pourquoi~?
\end{interviewq}

\begin{interviewq}[$\star\star\star$ \iqroles{trader, chercheur}]\label{iq:m1:shares-corporate-actions-indices:6}
Une société annonce une \omterm{def:m1:shares-corporate-actions-indices:rights}{augmentation de capital avec droits} 1-pour-3 avec une décote de 40~\%. Les actions
chutent de 12~\% à l'annonce. Est-ce «~à cause de la dilution~»~?
\end{interviewq}
